\section{The two imported analytic results}\label{sqt:sec:imports}
This section states exactly the consequences used from the primary sources. Their proofs are not reproduced here. The remaining arguments, including the reduction to their hypotheses, are given below.

\subsection{Absolute localization for geometric tables}
For $\theta,\phi\in(\R/2\pi\mathbb Z)^n$, put
\begin{equation}\label{sqt:eq:integrand}
\psi(z)=\frac{e^{-\ii z}}{2-e^{\ii z}},\qquad
\mathcal F(\theta,\phi)=\prod_{i,j=1}^n\psi(\theta_i+\phi_j).
\end{equation}
Fix a sufficiently small positive $\eta$. In Canfield--McKay's region $\mathcal R_\eta$, each angular vector lies in a semicircle, the centered row and column deviations have absolute value at most $\tfrac12n^{-1/2+\eta}$, and the circular sum $\mu$ of the two angular means obeys $|\mu|\leq\tfrac12n^{-1+4\eta}$. Means and deviations are evaluated in the corresponding small circular chart.

\begin{proposition}[Canfield--McKay, specialized]\label{sqt:prop:minor}
For all sufficiently large $n$,
\begin{equation}\label{sqt:eq:minor}
\int_{\mathcal R_\eta^c}|\mathcal F(\theta,\phi)|\,\dd\theta\,\dd\phi
 =O(e^{-c n^\eta})I_0,
\qquad I_0=\pi^{n-1/2}n^{1-n}e^{-1/4},
\end{equation}
where $c>0$ can be taken fixed.
\end{proposition}
This is Theorem~3 with the region in equation~(3.1) of the preprint \cite{CM}. At $m=n$, line sums $s=t=n$, and density $\lambda=1$, the left side of its condition~(1.1) is
\[
\frac{(1+2\lambda)^2}{4\lambda(1+\lambda)}
\left(1+\frac{5m}{6n}+\frac{5n}{6m}\right)=3.
\]
Thus the condition $3\leq a\log n$, with fixed $a,b>0$ and $a+b<1/2$, holds eventually. The polynomial upper bound on $\lambda$ in Theorem~3 is automatic. These are asymptotic conditions; no numerical starting index is inferred. We use the absolute minor-arc bound, not the coarser central Taylor remainder in the leading enumeration formula.

\subsection{A complex cumulant tail estimate}
Let $X=(X_1,\ldots,X_N)$ have independent coordinates taking values in a product domain $\Omega=\prod_{j=1}^N\Omega_j$. If $x,y\in\Omega$ and $V\subseteq[N]$, write $x\mathbin{\triangleright_V}y$ for the vector obtained by replacing the coordinates of $x$ in $V$ by those of $y$. For bounded measurable $f:\Omega\to\mathbb C$, define
\begin{align}
\Delta_V(f)&=\sup_{x,y\in\Omega}\left|
 \sum_{W\subseteq V}(-1)^{|W|}f(x\mathbin{\triangleright_W}y)\right|,\\
D_k(f)&=\max_{j\in[N]}\sum_{\substack{|V|=k\\j\in V}}\Delta_V(f).
\end{align}
Isaev's distribution-dependent quantities use
\[
\Delta_V(f,X)=\sup_{y\in\Omega}\left|\E\sum_{W\subseteq V}
 (-1)^{|W|}f(X\mathbin{\triangleright_W}y)\right|,
\]
and the corresponding sums $S_k(f,X)$. In particular, $S_k(f,X)\leq D_k(f)$. Cumulants of complex variables use ordinary products without complex conjugation.

\begin{proposition}[Isaev, sufficient form]\label{sqt:prop:complex}
Let $m$ be a positive integer. Suppose $f$ is bounded and measurable and, for every $1\leq t\leq m$,
\begin{equation}\label{sqt:eq:isaevhyp}
\sum_{k=t}^N e^{3\alpha k/2}2^t\binom{k-1}{t-1}D_k(f)
 \leq\alpha\leq\frac1{100}.
\end{equation}
Then there exists a complex $\delta$ such that
\begin{equation}\label{sqt:eq:isaev}
\E e^{f(X)}=(1+\delta)^N
 \exp\left(\sum_{r=1}^m\frac{\kappa_r(f(X))}{r!}\right),
\qquad |\delta|\leq e^{(100\alpha)^{m+1}}-1.
\end{equation}
\end{proposition}
This follows directly from Theorem~1.2 of \cite{Isaev} by $S_k\leq D_k$. The factor $(1+\delta)^N$ is important: at fixed $m$, if $N\alpha^{m+1}=o(1)$, the total multiplicative error is $1+O_m(N\alpha^{m+1})$. Bounds on mixed differences through order $m$ alone would not suffice; the sums in \eqref{sqt:eq:isaevhyp} continue to $N$.

The leading equivalent can alternatively be obtained from Barvinok--Hartigan's smooth-margin theorem \cite[Theorem~1.3]{BH}. Their absolute localization theorem is also compatible with the argument, but is not needed here. The binary graph-factor enumeration theorem of Isaev--McKay \cite{IM} is not used: its signless-Laplacian assumptions exclude a bipartite base graph. Nor does an enumeration theorem for Eulerian orientations \cite{IMZ} apply to the geometric entries. The distinction between Isaev's complex theorem and its earlier real-valued predecessor is essential because our cubic term is imaginary.
